\documentclass[ ../main.tex]{subfiles}
\providecommand{\mainx}{..}

% Permutation index can be a SECRET key in cipher systems.
% If wrong permutation index is used, can't really tell if it works or not.
% Given a single perf hash over A, maybe we can find multiple perfect hashes
% over a set of sets like A... another form of rate distortion.
%
%$\SetIndicator{\Set{A}} \colon \Set{X} \mapsto \BitSet$
%$\SetIndicator{\Set{A}}(x) \coloneqq \trunc(\hash(x),k) = \Fun{p}\Fun{?}[\Set{A}]$

\begin{document}
\section{Algebra of function composition}
If $\Fun{f} \colon \Set{X} \mapsto \Set{Y}$ and $\Fun{g} \colon \Set{Y} \mapsto \Set{Z}$ are injective functions, then $\Fun{g} \circ \Fun{f} \colon \Set{X} \mapsto \Set{Z}$ is an injective function.

A perfect hash function $\Fun{h}[\Set{A}] \colon \Set{X} \mapsto \Set{Y}$ is not injective.
However, its restriction,
\begin{equation}
	\Fun{h}[\Set{A}]\!\restriction_{\Set{A}} \colon \Set{A} \mapsto \Set{Y}
\end{equation}
is injective.
It immediately follows that $\Fun{g} \circ \Fun{h}[\Set{A}] \colon \Set{X} \mapsto \Set{Z}$ is injective over the restriction to $\Set{A}$, i.e., it is a perfect hash function over $\Set{A}$ of the type $\Set{X} \mapsto \Set{Z}$.

\begin{theorem}
Let $\Fun{g} \colon \Set{Y} \mapsto \Set{Z}$ be an injective function, $\Fun{h}[\Set{A}][r] \colon \Set{X} \mapsto \Set{Y}$ be a perfect hash function, and  $\Card{\Set{Z}} = (1 + \alpha) \Card{\Set{Y}}$, $\alpha \geq 0$.
The composition $\Fun{g} \circ \Fun{h}[\Set{A}][r] \colon \Set{X} \mapsto \Set{Z}$ is a perfect hash function $\Fun{h}[\Set{A}][r']$ where $r' = \frac{r}{1+\alpha}$.
\end{theorem}
\begin{proof}
It is given that $\Fun{h}[\Set{A}][r]$ has a load factor $r = \frac{\Card{\Set{A}}}{\Card{Y}}$, which means $\Card{\Set{Y}} = \frac{\Card{\Set{A}}}{r}$.
Since $\Fun{g} \circ \Fun{h}[\Set{A}][r]$ is an injective function into $\Set{Z}$, the load factor must be $r' = \frac{\Card{\Set{A}}}{\Card{\Set{Z}}}$.
It is also given that $\Card{\Set{Z}} = (1+\alpha) \Card{\Set{Y}}$.
Substituting $\Card{\Set{Z}}$ with $(1+\alpha) \Card{\Set{Y}}$, $\Card{\Set{Y}}$ with $\frac{\Card{\Set{A}}}{r}$, and simplifying yields the result $r' = \frac{r}{1+\alpha}$.
\end{proof}

\begin{theorem}
	Let $\Fun{g} \colon \Set{W} \mapsto \Set{X}$ be an injective function and let $\Fun{h}[\Set{A}][r] \colon \Set{X} 
	\mapsto \Set{Y}$ be a perfect hash function over $\Set{A}$.
	The composition $\Fun{h}[\Set{A}][r] \circ \Fun{g} \colon \Set{W} \mapsto \Set{Y}$, where $\Card{\Set{A}} = (1 + \alpha) \Card{\Set{W}}$, $\alpha \geq 0$, is a perfect hash function $\Fun{h}[\Set{W}][r'] \colon \Set{W} \mapsto \Set{Y}$ where $r' = \frac{r}{1+\alpha}$.
\end{theorem}
\begin{proof}
	It is given that $\Fun{h}[\Set{A}][r]$ has a load factor $r = \frac{\Card{\Set{A}}}{\Card{Y}}$, which means $\Card{\Set{Y}} = \frac{\Card{\Set{A}}}{r}$.
	It is also given that $\Card{\Set{Z}} = (1+\alpha) \Card{\Set{Y}}$.
	Since $\Fun{g} \circ \Fun{h}[\Set{A}][r]$ is an injective function into $\Set{Z}$, we see that the load factor must be $r' = \frac{\Card{\Set{A}}}{\Card{\Set{Z}}}$.
	However, we may substitute $\Card{\Set{Z}}$ with $(1+\alpha) \Card{\Set{Y}}$.
	Substituting $\Set{Y}$ with $\frac{m}{r}$ yields $\frac{\Card{\Set{A}}}{(1 + \alpha) \frac{\Card{\Set{A}}}{r}}$.
	Simplying yields the result $\frac{r}{1+\alpha}$.
\end{proof}

\begin{corollary}
	The composition $\Fun{h}[\Set{A}][r] \circ \Fun{g}[\Set{A}] \colon \Set{X} \mapsto \Set{Y}$ is a perfect hash function over some subset $\Set{T}$ of $\Set{W}$ such that $\Card{\Set{T}} \leq \Card{\Set{A}}$.
\end{corollary}

Suppose we have a perfect hash function $\Fun{h}[\Set{A}] \colon \Set{X} \mapsto \Set{Y}$.
Then, all of the perfect hash functions that may be generated by composing them with permutation functions are in the same \emph{equivalence class} with the property that every function in this equivalence contains the same distribution of collisions (and, of course, no collisions between elements of $\Set{A}$).

For instance, suppose we have a permutation function $\Fun{g} \colon \Set{Y} \mapsto \Set{Y}$ and a perfect hash function $\Fun{h}[\Set{A}] \colon \Set{X} \mapsto \Set{Y}$.
Then, the \emph{orbit} of $\Fun{g}^\infty \circ \Fun{h}[\Set{A}]$ generates a subset of this permutation equivalence class where $\Fun{g}^0 \coloneqq \Fun{id}[\Set{Y}]$.

Therefore, if we generate one perfect hash function over $\Set{A}$, we can immediately construct an entire family of perfect hash functions over $\Set{A}$.
This composition \emph{imbeds} the perfect hash function into a larger family of perfect hash functions.

If we limit our attention to bijective functions in $\Set{Y} \mapsto \Set{Y}$, then there are $\Card{\Set{Y}}!$ bijections in $\Set{Y} \mapsto \Set{Y}$ to choose from, i.e., generating one perfect hash function $\Fun{g}[\Set{A}] \colon \Set{X} \mapsto \Set{Y}$ immediately provides $\Card{\Set{Y}}!$ perfect hash functions of type $\Set{X} \mapsto \Set{Y}$ over $\Set{A}$ by \emph{permuting} the output.

We denote the family of perfect hash functions generated from a given perfect hash function $\Fun{f}[\Set{A}] \colon \Set{X} \mapsto \Set{Y}$ by $\sigma\left(\Fun{f}[\Set{A}]\right)$.
This is an \emph{equivalence class}, i.e., functions in this set are \emph{isomorphic}.
They have the property that for any pair of functions $\Fun{g},\Fun{h} \in \sigma \Fun{f}[\Set{A}]$, $\SetBuilder{ \Fun{g}(x) \in \Set{Y} }{ x \in \Set{B}} = $, where $\Set{B} \in \PS{\Set{X}}$



Using the previous construction of generating a family of perfect hash functions over $\Set{A}$ by permuting the output, the ratio of the number of such perfect hash functions to the total number of perfect hash functions is
\begin{equation}
	\frac
	{
		\left(\Card{\Set{Y}} - \Card{\Set{A}}\right)!
	}
	{
		\Card{\Set{Y}}^{\Card{\Set{X}} - \Card{\Set{A}}}
	}\,,
\end{equation}
where $\Set{A} \subseteq \Set{X}$, $\Card{\Set{X}} \geq \Card{\Set{Y}}$, and $\Card{\Set{A}} \leq \Card{\Set{Y}}$.



Suppose we have two perfect hash functions, $\Fun{g}[\Set{A}] \colon \Set{X} \mapsto \Set{Y}$ and $\Fun{f}[\Set{A}] \colon \Set{Y} \mapsto \Set{W}$.
Then, the function $\Fun{f} \circ \Fun{g} \colon \Set{W} \mapsto \Set{Y}$ is a perfect hash function over $\Set{A}$.








Suppose we have two perfect hash functions, $\Fun{g}[\Set{A}] \colon \Set{X} \mapsto \Set{Y}$ and $\Fun{f}[\Set{B}] \colon \Set{X} \mapsto \Set{Y}$.
Then, we may define the following higher-order operation.

We take the input $x \in \Set{X}$, apply it to both $\Fun{g}[\Set{A}]$ and $\Fun{f}[\Set{B}]$, yielding respectively $y_1$ and $y_2$ in $\Set{Y}$.
Next, we join the outputs in some way, e.g., the concatenation $y_1 y_2$.
Observe that no other $x \in \Set{A}$ hashes to $y_1$ and no $x \in \Set{B}$ hashes to $y_2$.
Therefore, no other $x \in \Set{A}$ hashes to $y_1 y_2$ and no other $x \in \Set{B}$ hashes to $y_1 y_2$.




\end{document}
